\cotefile{29}
\scanpage{119}
% archive sheet number: 118 (uncertain; alternative 119)

\handwritten{2}. Supposons donné une situation $(Y,V,H,V')$ analogue à $(X,\handwritten{\unsure{U}},G,U')$, et un morphisme $f:X\handwritten{\longrightarrow}Y$ tel que $f(U)\handwritten{\longrightarrow}f(V)$, un morphisme $f':U'\handwritten{\longrightarrow}V'$ compatible avec $U\handwritten{\longrightarrow}V$ induit par $f$, un morphisme de groupes $\handwritten{\phi}:G\handwritten{\longrightarrow}H$ tel que \deleted{$\unsure{f'(gx')=}$} $f'$ soit compatible avec les actions de $G$ et $H$ via $\handwritten{\phi}$, \deleted{\unsure{et}}

Le cas le plus courant sera celui où on se donne un point géométrique $a$ de $U$, \inserted{d'où} un point géométrique $b$ \inserted{$=f(a)$} de $V$, où $U'$ et $V'$ sont respectivement des quotients des revêtements universels de $U$ et $V$ basés en $a$ et $b$, $G$ et $H$ étant respectivement des quotients de $\handwritten{\pi}_{1}(U,a)$ et $\handwritten{\pi}_{1}(V,b)$, l'homomorphisme $f':U'\handwritten{\longrightarrow}V'$ étant induit par les homomorphismes canoniques sur les revêtements universels en $a$ resp. $b$\deleted{\unsure{;} \illegible{2} \unsure{ainsi} \unsure{donc} déterminé de}, \inserted{\deleted{ainsi que $\handwritten{\phi}$}} \deleted{façon unique $G$ et $H$ étant les \illegible{1}} $\handwritten{\phi}:G\handwritten{\longrightarrow}H$ étant induit par $f_{\handwritten{*}}:\handwritten{\pi}_{1}(U,a)\handwritten{\longrightarrow}\handwritten{\pi}_{1}(\unsure{B|V},b)$ ; de cette façon $f'$ et $\handwritten{\phi}$ sont déterminés sans ambiguïté par $f$, et leur existence signifie simplement que $f_{\handwritten{*}}:\handwritten{\pi}_{1}(U,a)\handwritten{\longrightarrow}\handwritten{\pi}_{1}(V,b)$ veut bien passer au quotient en un $\handwritten{\phi}:G\handwritten{\longrightarrow}H$ (i.e. applique un certain sous-groupe dans un certain sous-groupe). Considérons d'autre part le localisé strict $\handwritten{\overline}{Y}$ de $Y$ en le point géométrique $f(\handwritten{\bar}{x})$ de $Y$, d'où \deleted{\illegible{1}} $\handwritten{\bar}{V},\handwritten{\bar}{V}'$ comme dessus; supposons donné un prolongement $\handwritten{\overline}{Y'}$ de $\handwritten{\bar}{V}'$ en un schéma à groupe d'opérateurs $H$ sur $\handwritten{\bar}{Y}$, et supposons donné enfin un homomorphisme $\handwritten{\overline}{f'}:\handwritten{\overline}{X'}\handwritten{\longrightarrow}\handwritten{\overline}{Y'}$ qui prolonge l'homomorphisme \deleted{\unsure{ft|f'}} $\handwritten{\overline}{U'}\handwritten{\longrightarrow}\handwritten{\overline}{V'}$ induit par $f'$. Pour tout point $\handwritten{\unsure{\xi|\zeta}}$ de $\handwritten{\overline}{X'}_{\handwritten{\bar}{x}}$, $\handwritten{\overline}{f'}(\handwritten{\xi})$ est alors un point\deleted{\unsure{u|s}} de $\handwritten{\overline}{Y'}_{\handwritten{\bar}{y}}$, et on peut considérer le groupe d'inertie correspondant $H_{\handwritten{\overline}{f'}(\handwritten{\xi})}$. On a alors la relation évidente
\[
\handwritten{\phi}(G_{\handwritten{\xi}})\handwritten{\subset}H_{\handwritten{\overline}{f'}(\handwritten{\xi})}\ .
\tag*{\handwritten{(2.1)}}
\]
Dans le cas où le\deleted{\unsure{s}} groupe $G$ resp. $H$ opère transitivement sur $\handwritten{\overline}{X'}_{\handwritten{\bar}{x}}$ resp. $\handwritten{\overline}{Y'}_{\handwritten{\overline}{y'}}$, on peut donc dire que $\handwritten{\phi}$ \inserted{applique tout groupe \deleted{de}} \deleted{transforme la classe de conjugaison de sous-groupes} d'inertie de $G$ rel. à $x$ \inserted{\handwritten{\uline{dans}} un groupe d\deleted{e}\handwritten{'}} \deleted{\unsure{en|à} la classe de conj. des sous-groupes} d'inertie de $H$ rel à $y=f(x)$.

\deleted{D'autre part} Il reste alors à expliciter l'homomorphisme $\unsure{\handwritten{\overline}{f'}|\handwritten{\bar f}}:\handwritten{\overline{X'}\longrightarrow\overline{Y'}}$.

% Pencil enclosure around the paragraph beginning ``Le cas'', with a continuation arrow to the last paragraph.
\figure{\handwritten{\begin{tikzpicture}[x=0.5cm,y=0.5cm]
\draw (15,18) .. controls (14.6,18.2) and (14.8,18.5) .. (14,18.4) -- (1,18.4) .. controls (0.6,18.7) and (0.4,18) .. (0.4,17.6) -- (0.3,14) .. controls (0.8,12.5) and (0.5,11.4) .. (0.5,10.5) .. controls (1,10) and (1.3,10.2) .. (2,10.2) -- (7.4,10.2) .. controls (8.1,10.2) and (8.2,10.4) .. (8.2,10.9) -- (14.8,11.1);
\draw[->] (0.5,10.5) .. controls (-0.4,9) and (-0.3,3) .. (0.3,0.7) .. controls (0.5,0.1) and (0.9,0.5) .. (1.5,0.1);
\draw (-1.3,17.4) .. controls (-1.7,17.1) and (-1.7,16.9) .. (-1.5,16.6);
\end{tikzpicture}}}
